% Source: https://wiki.aalto.fi/display/Aaltothesis/Aalto+Thesis+LaTeX+Template

%%%%%%%%%%%%%%%%%%%%%%%%%%%%%%%%%%%%%%%%%%%%%%%%%%%%%%%%%%%%%%%%%%%%%%%%%%%%%%%%
%%%%%%%%%%%%%%%%%%%%%%%%%%%%%%%%%%%%%%%%%%%%%%%%%%%%%%%%%%%%%%%%%%%%%%%%%%%%%%%%
%%                                                                            %%
%% thesistemplate.tex version 4.10 (2025/06/30)                               %%
%% The LaTeX template file to be used with the aaltothesis.sty (version 4.10) %%
%% style file.                                                                %%
%% This package requires pdfx.sty v. 1.5.84 (2017/05/18) or newer.            %%
%%                                                                            %%
%% This is licensed under the terms of the MIT license below.                 %%
%%                                                                            %%
%% Written by Luis R.J. Costa.                                                %%
%% Currently developed at Teacher services, Learning Services of Aalto        %%
%% University by Luis R.J. Costa since May 2019.                              %%
%%                                                                            %%
%% Copyright 2017-2025 aaltothesis.cls by Luis R.J. Costa,                    %%
%% luis.costa@aalto.fi.                                                       %%
%% Copyright 2017-2018 Swedish translations in aaltothesis.cls by Elisabeth   %%
%% Nyberg and Henrik Wallén henrik.wallen@aalto.fi.                           %%
%% Copyright 2017-2018 Finnish documentation in the template opinnatepohja.tex%%
%% by Perttu Puska, perttu.puska@aalto.fi, and Luis R.J. Costa.               %%
%% Finnish documentation in the template opinnatepohja.tex translated from    %%
%% the English template documentation.                                        %%
%% Copyright 2025 English template thesistemplate.tex by Luis R.J. Costa,     %%
%% Maurice Forget, Henrik Wallén.                                             %%
%% Copyright 2018-2025 Swedish template kandidatarbetsbotten.tex by Henrik    %%
%% Wallen.                                                                    %%
%%                                                                            %%
%% Permission is hereby granted, free of charge, to any person obtaining a    %%
%% copy of this software and associated documentation files (the "Software"), %%
%% to deal in the Software without restriction, including without limitation  %%
%% the rights to use, copy, modify, merge, publish, distribute, sublicense,   %%
%% and/or sell copies of the Software, and to permit persons to whom the      %%
%% Software is furnished to do so, subject to the following conditions:       %%
%% The above copyright notice and this permission notice shall be included in %%
%% all copies or substantial portions of the Software.                        %%
%% THE SOFTWARE IS PROVIDED "AS IS", WITHOUT WARRANTY OF ANY KIND, EXPRESS OR %%
%% IMPLIED, INCLUDING BUT NOT LIMITED TO THE WARRANTIES OF MERCHANTABILITY,   %%
%% FITNESS FOR A PARTICULAR PURPOSE AND NONINFRINGEMENT. IN NO EVENT SHALL    %%
%% THE AUTHORS OR COPYRIGHT HOLDERS BE LIABLE FOR ANY CLAIM, DAMAGES OR OTHER %%
%% LIABILITY, WHETHER IN AN ACTION OF CONTRACT, TORT OR OTHERWISE, ARISING    %%
%% FROM, OUT OF OR IN CONNECTION WITH THE SOFTWARE OR THE USE OR OTHER        %%
%% DEALINGS IN THE SOFTWARE.                                                  %%
%%                                                                            %%
%%                                                                            %%
%%%%%%%%%%%%%%%%%%%%%%%%%%%%%%%%%%%%%%%%%%%%%%%%%%%%%%%%%%%%%%%%%%%%%%%%%%%%%%%%
%%                                                                            %%
%%                                                                            %%
%% An example for writting your thesis using LaTeX                            %%
%% Original version and development work by Luis Costa, changes to the text   %% 
%% in the Finnish template by Perttu Puska.                                   %%
%% Support for Swedish added 15092014                                         %%
%% PDF/A-b support added on 15092017                                          %%
%% PDF/A-2 support added on 24042018                                          %%
%% Layout design and typesettin changed 15072021                              %%
%%                                                                            %%
%% This example consists of the files                                         %%
%%       thesistemplate.tex (version 4.10) (for text in English)              %%
%%       opinnaytepohja.tex (version 4.10) (for text in Finnish)              %%
%%       kandidatarbetsbotten.tex (version 1.20) (for text in Swedish)        %%
%%       thesistemplate_short.tex (version 4.10) (abridged for text in        %%
%%                                                English)                    %%
%%       aaltothesis.cls (version 4.10)                                       %%
%%       linediagram.pdf (graphics file)                                      %%
%%       curves.pdf      (graphics file)                                      %%
%%       ledspole.jpg    (graphics file)                                      %%
%%                                                                            %%
%%                                                                            %%
%% Typeset in Linux with                                                      %%
%% pdflatex: (recommended method)                                             %%
%%             $ pdflatex thesistemplate                                      %%
%%             $ pdflatex thesistemplate                                      %%
%%                                                                            %%
%%   The result is the file thesistemplate.pdf that is PDF/A compliant, if    %%
%%   you have chosen the proper \documenclass options (see comments below)    %%
%%   and your included graphics files have no problems.                       %%
%%                                                                            %%
%%                                                                            %%
%% Explanatory comments in this example begin with the characters %%, and     %%
%% changes that the user can make with the character %                        %%
%%                                                                            %%
%%%%%%%%%%%%%%%%%%%%%%%%%%%%%%%%%%%%%%%%%%%%%%%%%%%%%%%%%%%%%%%%%%%%%%%%%%%%%%%%
%%%%%%%%%%%%%%%%%%%%%%%%%%%%%%%%%%%%%%%%%%%%%%%%%%%%%%%%%%%%%%%%%%%%%%%%%%%%%%%%
%%
%% WHAT is PDF/A
%%
%% PDF/A is the ISO-standardized version of the pdf. The standard's goal is to
%% ensure that he file is reproducable even after a long time. PDF/A differs
%% from pdf in that it allows only those pdf features that support long-term
%% archiving of a file. For example, PDF/A requires that all used fonts are
%% embedded in the file, whereas a normal pdf can contain only a link to the
%% fonts in the system of the reader of the file. PDF/A also requires, among
%% other things, data on colour definition and the encryption used.
%% Currently three PDF/A standards exist:
%% PDF/A-1: based on PDF 1.4, standard ISO19005-1, published in 2005.
%%          Includes all the requirements essential for long-term archiving.
%% PDF/A-2: based on PDF 1.7, standard ISO19005-2, published in 2011.
%%          In addition to the above, it supports embedding of OpenType fonts,
%%          transparency in the colour definition and digital signatures.
%% PDF/A-3: based on PDF 1.7, standard ISO19005-3, published in 2012.
%%          Differs from the above only in that it allows embedding of files in
%%          any format (e.g., xml, csv, cad, spreadsheet or wordprocessing
%%          formats) into the pdf file.
%% PDF/A-4: based on PDF 2.0, standard ISO19005-4, published in November 2020.
%%
%% PDF/A-1 files are not necessarily PDF/A-2 -compatible and PDF/A-2 are not
%% necessarily PDF/A-1 -compatible.
%% Standards PDF/A-1, PDF/A-2 and PDF/A-3 have two levels:
%% b: (basic) requires that the visual appearance of the document is reliably
%%    reproduceable.
%% a (accessible) in addition to the b-level requirements, specifies how
%%   accessible the pdf file is to assistive software, say, for the physically
%%   impaired.
%% The PDF/A-4 standard does not have additional levels like in the earlier
%% standards.
%% For more details on PDF/A, see, e.g., 
%% https://www.loc.gov/preservation/digital/formats/fdd/fdd000318.shtml or
%% https://www.pdfa.org/resource/iso-19005-pdfa/
%%
%%
%% WHICH PDF/A standard should my thesis conform to?
%%
%% Either to the PDF/A-2b or the PDF/A-1b standard. If all the figures and
%% graphs used in thesis work do not require transparency features, use either
%% PDF/A-1b or PFDF/A-2b. If you have figures with transparency
%% characteristics, use the PDF/A-2b standard. However, drawing applications
%% often use the transparency parameter, setting it to zero, to specify opacity
%% and get the basic 2-D visualisation. As a result, validation of PDF/A-1b
%% will fail. Use PDF/A-2b if PDF/A-1b validation fails.
%% Do not use the PDF/A-3b standard for your thesis.
%% The font to be used are specified in this templatenand they should not be
%% changed. In addition to not adhering to Aalto's visual guidelines, you may
%% have difficulties in producing a PDF/A-compliant PDF.
%%
%%
%% Validate your PDF/A file at https://www.pdf-online.com/osa/validate.aspx
%%
%%
%% WHAT graphics format can I use to produce my PDF/A compliant file?
%%
%% When using pdflatex to compile your work, favour the use of pdf, but you can
%% use the jpg or png format especially for photographs. You will have PDF/A 
%% compliance problems with figures in pdf if the fonts are not embedded in the
%% pdf file.
%% If you choose to use latex to compile your work, the only acceptable file
%% format for your figure is eps. DO NOT use the ps format for your figures.
%%
%% USE one of the following three \documentclass set-ups:
%% * the first when using pdflatex to directly typeset your document in the
%%   chosen pdf/a format for online publishing (centred page layout),
%% * the second for one-sided printing your thesis with the layout (wide left 
%%   margin), or
%% * the third for two-sided printing.
%%
\documentclass[english, 12pt, a4paper, elec, utf8, a-2b, online]{aaltothesis}
%\documentclass[english, 12pt, a4paper, elec, utf8, a-2b, print]{aaltothesis}
%\documentclass[english, 12pt, a4paper, elec, utf8, a-2b, print, twoside]{aaltothesis}

%% Use the following options in the \documentclass macro above:
%% your school: arts, biz, chem, elec, eng, sci
%% the character encoding scheme used by your editor: utf8, latin1
%% thesis language: english, finnish, swedish
%% make an archiveable PDF/A-1b or PDF/A-2b compliant file: a-1b, a-2b
%%                    (with pdflatex, a normal pdf containing metadata is
%%                     produced without the a-*b option)
%% typset for online document or print on paper: online, print
%%        online: typeset in symmetric layout and blue hypertext for online
%%                publishing
%%        print: typeset in a symmetric layout and black hypertext for printing
%%               on paper
%%          two-side printing: twoside (default is one-sided printing)
%%               typeset in a wide margin on the binding side of the page and
%%               black hypertext. Use with print only.
%%

%% Use one of these if you write in Finnish (or use the Finnish template
%% opinnaytepohja.tex)
%\documentclass[finnish, 12pt, a4paper, elec, utf8, a-1b, online]{aaltothesis}
%\documentclass[finnish, 12pt, a4paper, elec, utf8, a-1b, print]{aaltothesis}
%\documentclass[finnish, 12pt, a4paper, elec, utf8, a-1b, print, twoside]{aaltothesis}

%% Use one of these if you write in Swedish (or use the Swedish template
%% kandidatarbetsbotten.tex)
%\documentclass[swedish, 12pt, a4paper, elec, utf8, a-2b, online]{aaltothesis}
%\documentclass[swedish, 12pt, a4paper, elec, utf8, a-2b]{aaltothesis}
%\documentclass[swedish, 12pt, a4paper, elec, dvips, online]{aaltothesis}

%% FOR USERS OF AMS PACKAGES:
%% * newtxmath used in this template loads amsmath, so
%%   you needn't load it. If you want to use options in amsmath, load it here, 
%%   before \setupthesisfonts below to pass the options to amsmath.
%% * If you want to use amsthm, load it here before \setupthesisfonts to avoid
%%   a clash with newtxmath.
%% * If using amsmath with options and you want to use amsthm, load amsthms
%%   after amsmath, as described in the amsthm documentation.
%% * Don't use amsbsym or amsfonts. The symbols [and macros] there are defined in
%%   newtxmath and so clash if used.
%\usepackage[options]{amsmath}
%\usepackage{amsthm}

%% DO NOT MOVE OR REMOVE \setupthesisfonts
\setupthesisfonts

%%
%% Add here the packages you need
%%
\usepackage{graphicx}
\usepackage{hyperref}
\usepackage{tabularx}
\usepackage{xurl}
\usepackage[numbers]{natbib}

\usepackage{todonotes}
\setuptodonotes{inline}

% Rewrite \autoref
\addto\extrasenglish{
  \def\sectionautorefname{Section}
  \def\subsectionautorefname{Subsection}
}

%% For tables that span multiple pages; used to split a paraphrasing example in
%% the appendix. If you don't need it, remove it.
\usepackage{longtable}

%% A package for generating Creative Commons copyright terms. If you don't use
%% the CC copyright terms, remove it, since otherwise undesired information may
%% be added to this document's metadata.
\usepackage[type={CC}, modifier={by-nc-sa}, version={4.0}]{doclicense}
%% Find below three examples for typesetting the CC license notice.


%% Edit to conform to your degree programme
%% Capitalise the words in the name of the degree programme: it's a name
\degreeprogram{Computer, Communication and Information Sciences}
%%

%% Your major
%%
\major{Computer Science}
%%

%% Choose one of the three below
%%
%\univdegree{BSc}
\univdegree{MSc}
%\univdegree{Lic}
%%

%% Your name (self explanatory...)
%%
\thesisauthor{Minghao Guo}
%%

%% Your thesis title and possible subtitle comes here and possibly, again,
%% together with the Finnish or Swedish abstract. Do not hyphenate the title
%% (and subtitle), and avoid writing too long a title. Should LaTeX typeset a
%% long title (and/or subtitle) unsatisfactorily, you might have to force a
%% linebreak using the \\ control characters. In this case...
%% * Remember, the title should not be hyphenated!
%% * A possible 'and' in the title should not be the last word in the line; it
%%   begins the next line.
%% * Specify the title (and/or subtitle) again without the linebreak characters
%%   in the optional argument in box brackets. This is done because the title
%%   is part of the metadata in the pdf/a file, and the metadata cannot contain
%%   linebreaks.
%%
\thesistitle{Hypermedia-Constrained Generative UI}
%\thesistitle[Title of the thesis]{Title of\\ the thesis}
%%
%% Either remove or leave \thesissubtitle{} empty if you don't use it
%%
\thesissubtitle{A Deterministic Framework for Human-Agent Interaction}
%\thesissubtitle[Subtitle of the thesis]{Subtitle of\\ the thesis}
%\thesissubtitle{}

%%
\place{Espoo}
%%

%% The date for the bachelor's thesis is the day it is presented
%%
\date{30 October 2026}
%%

%% Thesis supervisor
%% Note the "\" character in the title after the period and before the space
%% and the following character string.
%% This is because the period is not the end of a sentence after which a
%% slightly longer space follows, but what is desired is a regular interword
%% space.
%%
\supervisor{Prof.\ Petri Vuorimaa}
%%

%% Advisor(s)---two at the most---of the thesis. Check with your supervisor how
%% many official advisors you can have.
%%
\advisor{University Teacher Juho Vepsäläinen (DSc)}
%\advisor{Ms Elsa Expert (MSc)}
%%

%% If you do your thesis work in a company of other institute, give the name of
%% the company or instution here. Otherwise, leave the macro empty, comment it
%% out, or remove it. This will remove this field from the abstract page.
%%
% \collaborativepartner{Company or institute name (if relevant)}
%%

%% Aaltologo: syntax:
%% \uselogo{?|!|'|aalto?|aalto!|aalto'|<empty>}
%% The logo language is set to be the same as the thesis language.
%%
%\uselogo{?}
%\uselogo{!}
\uselogo{}
%\uselogo{aalto?}
%\uselogo{aalto!}
%\uselogo{aalto'}
%\uselogo{}
%%



%%%%%%%%%%%%%%%%%%               COPYRIGHT TEXT               %%%%%%%%%%%%%%%%%%
%%%%%%%%%%%%%%%%%%%%%%%%%%%%%%%%%%%%%%%%%%%%%%%%%%%%%%%%%%%%%%%%%%%%%%%%%%%%%%%%

%% Copyright of a work is with the creator/author of the work regardless of
%% whether the copyright mark is explicitly in the work or not. You may, if you
%% wish---we encourage you to do so---publish your work under a Creative
%% Commons license (see creativecommons.org), in which case the license text
%% must be visible in the work. Write here the copyright text you want using the
%% macro \copyrighttext, which writes the text into the metadata of the pdf file
%% as well.
%%
%% Syntax:
%% \copyrighttext{metadata text}{text visible on the page}
%%
%% CHOOSE ONE OF THE COPYRIGHT NOTICE STYLES BELOW.
%% IF USING THE CC TERMS, CHOOSE THE LICENSE YOU WANT TO USE.
%% The different CC licenses are listed at 
%% https://creativecommons.org/about/cclicenses/.
%% If you use the icons from the doclicense.sty package, add the package above
%% (\usepackage{doclicense}).
%% IMPORTANT NOTE!! Manually write the CC text in the \copyrighttext metadata
%% text field.
%%
%% NOTE: In the macros below, the text written in the metadata must have a
%% \noexpand macro before the \copyright special character. When not in pdf/a
%% mode (i.e. a-1b or a-2b are not specified in \documentclass), two \noexpands
%% are required in the metadata text to correctly render the copyright mark in
%% the pdf metadata. In pdf/a mode one \noexpand suffices.
%%
%% EXAMPLE OF PLAIN COPYRIGHT TEXT
%% The macros \copyright and \year below must be separated by the \ character 
%% (space chacter) from the text that follows. The macros in the argument of the
%% \copyrighttext macro automatically insert the year and the author's name.
%% (Note! \ThesisAuthor is an internal macro of the aaltothesis.cls class file).
%%
%\copyrighttext{Copyright \noexpand\textcopyright\ \number\year\ \ThesisAuthor}
%{Copyright \textcopyright{} \number\year{} \ThesisAuthor}
%%
%% Of course, the same text could have simply been written as
%% \copyrighttext{Copyright \noexpand\copyright\ 2018 Eddie Engineer}
%% {Copyright \copyright{} 2022 Eddie Engineer}
%%
%% EXAMPLES OF CC LICENSE: different ways to display the same license
%% 1. A simple Creative Commons license text with a link to the copyright notice:
%\copyrighttext{\noexpand\textcopyright\ \number\year. This work is 
%	licensed under a CC BY-NC-SA 4.0 license.}{\textcopyright{} 
%	\number\year. This work is licensed under a 
%	\href{https://creativecommons.org/licenses/by-nc-nd/4.0/}{CC BY-NC-SA 4.0} 
%	license.}
%
%% To get the URL of the license of your choice, go to 
%% https://creativecommons.org/about/cclicenses/, click on the chosen license
%% you want to use, and copy-and-paste the URL in the macro \href above.
%%
%% 2. A short Creative Commons license text containing the respective CC icons
%% (requires the package doclicense.sty to be added in the preamble as done
%% above) and a link to the corresponding Creative Commons license webpage (see
%% the doclicense package documentation for other license icons):
%\copyrighttext{\noexpand\textcopyright\ \number\year. This work is licensed
%	under a CC BY-NC-SA 4.0 license.}{
%	\parbox{95mm}{\noindent\textcopyright\ \number\year. \doclicenseText} 
%	\hspace{1em}\parbox{35mm}{\doclicenseImage}
%}
%%
%% 3. An expanded Creative Commons license text containing the respective CC
%% icons text and as generated by the doclicense.sty package (the license is set
%% via package options in \usepackage[options]{doclicense} above; see the
%% doclicense package documentation for other license texts and icons):
\copyrighttext{\noexpand\textcopyright\ \number\year. This work is 
	licensed under a Creative Commons "Attribution-NonCommercial-ShareAlike 4.0 
	International" (BY-NC-SA 4.0) license.}{\noindent\textcopyright\ \number
	\year \ \doclicenseThis}
%%%%%%%%%%%%%%%%%%%%%%%%%%%%%%%%%%%%%%%%%%%%%%%%%%%%%%%%%%%%%%%%%%%%%%%%%%%%%%%%

%% You can prevent LaTeX from writing into the xmpdata file (it contains all the 
%% metadata to be written into the pdf file) by setting the writexmpdata switch
%% to 'false'. This allows you to write the metadata in the correct format
%% directly into the file thesistemplate.xmpdata.
%\setboolean{writexmpdatafile}{false}


%% All that is printed on paper starts here
%%
\begin{document}

%% Create the coverpage
%%
\makecoverpage

%% Typeset the copyright text.
%% If you wish, you may leave out the copyright text from the human-readable
%% page of the pdf file. This may seem like a attractive idea for the printed
%% document especially if "Copyright (c) yyyy Eddie Engineer" is the only text
%% on the page. However, the recommendation is to print this copyright text.
%%
\makecopyrightpage

\clearpage
%% Note that when writing your thesis in English, place the English abstract
%% first followed by the possible Finnish or Swedish abstract.

%% Abstract text
%% All the details (name, title, etc.) on the abstract page appear as specified
%% above. Add your abstarct text with paragraphs here to have paragraphs in the
%% visible abstract page. Nonetheless, write the abstarct text without
%% paragraphs in the macro \thesismacro so that it is added to the metadata as
%% well.
%%
\begin{abstractpage}[english]

\todo{Write the abstract after everything settled.}

  % The abstract is a short description of the essential contents of the thesis,
  % usually in one paragraph: what was studied and how and what were the main
  % findings.

  % For a Finnish thesis, the abstract should be written in both Finnish and
  % English; for a Swedish thesis, in Swedish and English. The abstracts for
  % English theses written by Finnish or Swedish speakers should be written in
  % English and either in Finnish or in Swedish, depending on the student’s
  % language of basic education. Students educated in languages other than Finnish
  % or Swedish write the abstract only in English. Students may include a second
  % or third abstract in their native language, if they wish.

  % The abstract text of this thesis is written on the readable abstract page as
  % well as into the pdf file's metadata via the \verb+\thesisabstract+ macro
  % (see comment in this \TeX{} file above). Write here the text that goes onto
  % the readable abstract page. You can have special characters, linebreaks, and
  % paragraphs here. Otherwise, this abstract text must be identical to the
  % metadata abstract text.
  
  % If your abstract does not contain special characters and it does not require
  % paragraphs, you may take advantage of the \verb+\abstracttext+ macro (see the
  % comment in this \TeX{} file below).
\end{abstractpage}

%% The text in the \thesisabstract macro is stored in the macro \abstractext, so
%% you can use the text metadata abstract directly as follows:
%%
%\begin{abstractpage}[english]
%	\abstracttext{}
%\end{abstractpage}

%% Force a new page so that the possible Finnish or Swedish abstract does not
%% begin on the same page
%%
% \newpage
%%
%% Abstract in Finnish. Delete if you don't need it. 
%%
%% Respecify those fields that differ from the earlier specification or simply
%% respecify all fields.
% \thesistitle{Opinnäyteen otsikko}
% \thesissubtitle{Opinnäytteen mahdollinen alaotsikko}
% \supervisor{Prof.\ Pirjo Professori}
% \advisor{Yliopisto-opettaja Juho Vepsäläinen}
% \degreeprogram{Elektroniikka ja sähkötekniikka}
% \major{Sopiva pääaine}
% \collaborativepartner{Yhtiön tai laitoksen nimi (tarvittaessa)}
% \date{30.6.2025}
% %% The keywords need not be separated by \spc now.
% \keywords{Vastus, resistanssi, lämpötila}
% %% Abstract text
% \begin{abstractpage}[finnish]
% Tiivistelmä on lyhyt kuvaus työn keskeisestä sisällöstä usein yhtenä kappaleena:
% mitä tutkittiin ja miten sekä mitkä olivat tärkeimmät tulokset. Suomenkielisen
% opinnäytteen tiivistelmä kirjoitetaan suomeksi ja englanniksi ja ruotsinkielisen
% vastaavasti ruotsiksi ja englanniksi. Suomen- tai ruotsinkielisten
% opiskelijoiden, joiden opinnäytteen kieli on englanti, tulee kirjoittaa
% tiivistelmänsä englanniksi ja koulusivistyskielellään. Muiden kuin
% koulusivistyskieleltään suomen- tai ruotsinkielisten tulee kirjoittaa
% tiivistelmänsä vain englanniksi. Opiskelija voi halutessaan lisätä
% opinnäytteeseensä toisen tai kolmannen tiivistelmän omalla äidinkielellään.
% Tämän opinnäytteen tiivistelmäteksti kirjoitetaan opinnäytteen luettavan osan
% lomakkeen lisäksi myös pdf-tiedoston metadataan. Kirjoita tähän metadataan
% kirjoitettavaa teksti. Metadatatekstissa ei saa olla erikoismerkkejä,
% rivinvaiho- tai kappaleenjakomerkkiä, joten näitä merkkeja ei saa käyttää tässä.
% Jos tiivistelmäsi ei sisällä erikoimerkkejä eikä kaipaa kappaleenjakoa, voit
% hyödynttää makroa abstracttext luodessasi lomakkeen tiivistelmää (katso
% kommentti tässä TeX-tiedostossa alla). Metadatatiivistelmatekstin on muuten 
% oltava sama kuin lomakkeessa oleva teksti.

% \end{abstractpage}

% %% Force new page so that the Swedish abstract starts from a new page
% \newpage

% %% Swedish abstract. Delete it if you don't need it. 
% %% 
% %% Respecify those fields that differ from the earlier specification or simply
% %% respecify all fields.
% \thesistitle{Arbetets titel}
% \supervisor{Prof.\ Pirjo Professori}
% \advisor{TkD Alan Advisor} %
% \advisor{DI Elsa Expert}
% \degreeprogram{Electronik och electroteknik}
% \collaborativepartner{Company or institute name in Swedish (if relevant)}
% %\date{30.6.2025}
% %% Abstract keywords
% \keywords{Nyckelord p\aa{} svenska, temperatur}
% %% Abstract text
% \begin{abstractpage}[swedish]
% Sammandraget är en kort beskrivning av arbetets centrala innehåll: vad 
% undersöktes, hur undersöktes det och vilka var de viktigaste resultaten?

% I lärdomsprov som skrivs på svenska skrivs sammandraget på svenska och engelska, 
% på motsvarande sätt skrivs sammandraget på finska och engelska i lärdomsprov på 
% finska. Finsk- eller svenskspråkiga studerande som skriver sitt lärdomsprov på 
% engelska ska skriva sammandraget på engelska och på sitt skolutbildningsspråk. 
% Studerande vars skolutbildningsspråk inte är svenska eller finska skriver 
% sammandraget endast på engelska. Den studerande kan om hen så önskar lägga till 
% ett andra eller tredje sammandrag på sitt eget modersmål. Sammandraget fungerar 
% då ofta som mognadsprov och bör i så fall vara minst 300 ord långt. Information 
% om mognadsprov på svenska finns på MyCourses:\\
% \url{https://mycourses.aalto.fi/course/view.php?id=26872}.
% \end{abstractpage}


\dothesispagenumbering{}

% %% Preface
% %%
% %% This section is optional. Remove it if you do not want a preface.
% \mysection{Preface}
% %\mysection{Esipuhe}
% I want to thank Professor Pirjo Professor and my instructors Dr Alan Advisor and
% Ms Elsa Expert for their guidance.

% I also want to thank my partner for keeping me sane and alive.

% \vspace{5cm}
% Otaniemi, 30 June 2025\\

% \vspace{5mm}
% {\hfill Eddie E.\ Engineer \hspace{1cm}}

% %% Force a new page after the preface
% %%
% \newpage

%% Table of contents. 
%%
\thesistableofcontents

\input{_abbreviations}

%% \clearpage is similar to \newpage, but it also flushes the floats (figures
%% and tables).
%%
\cleardoublepage

%% Text body begins.
%%
\section{Introduction}
\label{sec:introduction}

%% Leave page number of the first page empty
%% 
\thispagestyle{empty}

% 1. Describe the overall context + bring several background references
% 2. Show the research gap (why to do this particular research)
% 3. Bring up your research question(s)
% 4. Describe the structure of the work + mention research method

% Research methods: literature review, design science prototype, action-validity evaluation

\todo{A paragraph filled with ancient web development and the modern LLM revolution, leads to topic}

\subsection{Why this topic}

\todo{Find supporting materials for this and make each bullet a point}
\begin{itemize}
    \item Cognitive load and inefficiency of pure conversational interfaces in complex, state-dependent web tasks.
    \item Current Agentic/Generative UI systems often have to infer available actions from unstructured context, human-facing pages, or weakly structured APIs.
    \item Utilizing HTML hypermedia representations that expose state-dependent affordances as a trusted information layer for state-aware agentic UI generation.
\end{itemize}

%\cite{jarvi_algorithms_2009}
%\cite{yang2025agentic}
%\citet{jarvi_algorithms_2009}

\subsection{Research Questions}

\begin{enumerate}
    \item[] \textbf{RQ1:} What are the fundamental limitations of CUIs that necessitate Agentic UIs, and how do state-of-the-art Generative UI frameworks differ in terms of their benefits and drawbacks?
    
    % \item[] \textbf{RQ2:} How can a system architecture be designed to parse backend hypermedia constraints (application state and affordances) and deterministically map them to frontend components based on user intent?
    \item[] \textbf{RQ2:}  How can HTML hypermedia affordances, such as links, forms, actions, and element states, be extracted and represented as a trusted constraint layer for state-aware agentic UI generation?
    

    % \item[] \textbf{RQ3:} Measured via computational HCI models (e.g., Keystroke-Level Model), to what extent do hypermedia-generated GUIs reduce interaction complexity and operation step counts for complex web tasks compared to standard conversational interfaces?
    \item[] \textbf{RQ3:} In a simulated state-dependent workflow, to what extent does an HTML hypermedia-constrained agentic UI reduce invalid action suggestions compared with an unconstrained or weakly structured baseline?
\end{enumerate}

\todo{Point out the contribution of this work: a framework comparison, an HTML hypermedia-constrained architecture, and a prototype evaluation focused on action validity rather than general UI quality.}


\subsection{Structure of the work}

\todo{Describe the structure of the work and tie research questions to it. 

Better do this at the final stage}


\clearpage

\section{Background}
\label{sec:background}
% May change title to Background or more specific meaning

Software agents built on large language models can now operate the web on a person's behalf, and the interfaces through which people direct them are being redesigned around that fact. This chapter discusses how agents increasingly meet their users through generated interfaces. Existing frameworks for such interfaces are reviewed, the question of what grounds the actions they offer is raised, and HTML hypermedia is proposed as a solution.

\subsection{Internet, Human, and Machine}\label{sec:bg:internet}

Web interaction evolved from pages authored for human reading toward pages that non-human clients also need to act on. The earliest such clients like web crawlers only read while later ones increasingly act: scripted automation at first, and most recently LLM-based agents that follow links, fill forms, and complete transactions on the user's behalf. \todo{cite: general web history / rise of non-human web clients}

\todo{figure: human-facing web pages, machine-facing web pages... }

Because pages are written for humans, a machine client has no direct access to application meaning: it must recover what a page is, and what can currently be done with it, from artifacts intended for human perception --- markup structure, layout, and visible text. 
This recovery is indirect and brittle, and the difficulty motivated a long line of work on giving the web machine-readable semantics, most prominently the Semantic Web programme and its lighter-weight descendants such as embedded structured annotations. \todo{cite: semantic web, schema.org, microdata, RDFa, JSON-LD, etc.}
% https://arxiv.org/html/2507.10644v1#bib.bib87 may be useful for the semantic-web citation
(Maybe: shall we extend a little bit to make a list?)

The semantic-web approach asked authors to maintain a second, machine-facing representation of their application in parallel with the human-facing page. Annotations that are not needed to render the page for people are exercised by nothing the author depends on, and so they are the first thing to fall out of date. \todo{cite: find supporting old papers}

\subsection{From Conversational Agents to Generative UI}\label{sec:bg:genui}

Large language models entered this picture from the opposite direction: rather than making pages more readable to machines, they made machines better readers of artifacts meant for humans. Their first and still dominant form is conversational. Chat interfaces in the ChatGPT style proved remarkably capable for open-ended language tasks, and they inherit a long lineage of conversational-interface research \cite{mctear2017rise}.

These models were then given the ability to act with scaffolding for planning, tool use, and memory, they can call tools, operate websites, and carry out multi-step tasks against external systems, and an ecosystem of such agents has grown around the web in particular \cite{yang2025agentic, ning2025survey, zhou2024webarena, gur2023realworld, song2025browsing}. This turned how a human should interact with a capable agent into a question.

For these richer, structured, stateful tasks, pure conversation shows limits:
\begin{enumerate}
    \item \textbf{Lossy communication:} natural language is a lossy, serial channel for what are really structured choices: selecting one of many flights, configuring a filter, or filling a constrained form is painful to conduct sentence by sentence.
    \item \textbf{Hidden action space:} narrating actions hides the actual action space from the human, who cannot see what the agent could have done, verify what it actually did, or take over precisely \cite{zhang2025exploring}.
\end{enumerate}

Hence generative UI: instead of only talking, the agent renders task-specific interaction surfaces --- controls, forms, choices --- fitted to the current step of the task \cite{leviathan2025generative, aisdkGUI}. The industry is already moving this way: mainstream assistants increasingly wrap their responses in generated widgets and forms even outside web-operation tasks, which corroborates that pure text is insufficient as an interaction surface.
\todo{cite: Gemini / GPT generative-UI product direction, maybe webpage} 

This thesis narrows to the demanding end of that: generated controls that correspond to real actions on an external, stateful application. Namely the correctness of the surface --- does this control exist, and is it valid now? The rest of this chapter is organized around that problem: a generated surface can present controls that are invalid, stale, or outright hallucinated relative to what the application can actually do at this moment. 
\todo{cite: LLM hallucination}

\subsection{Generative UI and Agentic UI Frameworks}\label{sec:bg:frameworks}

\todo{figure: human, agent, application data-flow pipeline (intent, action selection, execution, new state, rendering); locate each framework on it}

To compare existing systems... /todo{describe the figure}

Two concerns run through this loop and are worth separating, because existing systems tend to address only one of them. The rendering side concerns how agent output is drawn for the human. The action-source side concerns where the set of actions the agent may take, and may offer, comes from. The two are independent, and the rest of this section groups systems by the second.

\paragraph{Rendering-side}

AG-UI defines an event and transport protocol between agents and front ends \cite{aguiOverview}, Vercel's AI SDK maps model output to registered interface components \cite{aisdkGUI}, and CopilotKit embeds agent interaction into existing applications. \todo{cite: CopilotKit} These systems are complementary to this thesis: they answer how a generated surface is drawn, not what it may legitimately contain, and they could sit on top of the framework developed here.
\todo{expand this}

\paragraph{Source-side}

The weakest grounding is inference from the live page. Browser and computer-use agents recover their action set from the DOM, the accessibility tree, or screenshots, and realistic web tasks remain difficult and brittle when the agent must guess in this way what a page affords \cite{zhou2024webarena, song2025browsing}. 
% This is the indirect-recovery problem of Section~\ref{sec:bg:internet} replayed with a stronger model. 
Magentic-UI is the generative-UI exemplar ... \cite{mozannar2025magenticui}.

A stronger answer registers the actions on the client at build time: the developer declares a set of tools or components that the model may invoke, explicit and typed \cite{aisdkGUI}. This grounds the vocabulary but not the state. The registry is static and application-specific, and nothing prevents it from drifting away from what the backend currently allows.

The strongest existing answer is a server-declared contract, and its current form is the Model Context Protocol (MCP). MCP deserves full credit on this axis: the contract is machine-readable, owned by the server rather than assumed by the client, gives every tool a typed input schema, and can update the tool list at runtime.
\todo{cite: MCP}

\todo{table: five-dimension framework comparison; rows: DOM-inference agents, build-time registries, MCP, hypermedia affordances (this thesis, forward reference to the Design chapter); cite each row's primary system}

\subsection{HTML Hypermedia, State, and Affordances}\label{sec:bg:hypermedia}

\subsubsection{REST and the hypermedia constraint}

Hypermedia denotes representations that carry both content and the controls for moving to other application states \cite{hypermediaBook}. In HTML those controls are links, forms, and buttons, and in server-driven variants such as htmx, attributes that declare further exchanges with the server \cite{htmxlDocs}. Hypermedia formats aimed only at machines also exist, such as the JSON-based HAL and Siren, but HTML is the one that matters here, for a reason this section builds toward: it is the only widely deployed hypermedia format whose representations are at the same time a human interface.

This is not a new idea but the original architecture of the web. Fielding's REST names the principle hypermedia as the engine of application state: transitions should be driven by controls that the server places in the current representation, not by out-of-band contracts the client holds separately. Read against Section~\ref{sec:bg:frameworks}, this identifies the pattern that MCP repeats. WSDL and OpenAPI were the out-of-band contracts of their eras; a tool catalog obtained at connection time is the same client-held description, separate from the representation the application actually serves.
\todo{cite: REST, Hypermedia} 

The controls in a representation can be read as its affordances: the state-dependent set of actions the application currently makes available \cite{gidey2025affordance}. The term is used here in the hypermedia community's sense \cite{hypermediaBook}. 

\subsubsection{Consequences}

The observation this thesis is built on follows. In a hypermedia application, the HTML the server returns is one artifact playing three roles at once: the interface the human sees, the action space an agent can parse, and the set of controls through which the application is actually operated. These three cannot drift apart, because they are not three things. An agent that takes those controls as its source of valid actions is not guessing from arbitrary DOM structure, and not trusting a separately maintained description; it is reading the same contract the human user is given. A separate contract such as MCP cannot offer this property, however well specified, because separateness is exactly what the property excludes.

\begin{enumerate}
    \item \textbf{No drift:} the affordance description is a byproduct of building the application, not a second artifact to maintain, so the incentive failure of Section~\ref{sec:bg:internet} does not arise.
    \item \textbf{Native state-dependence:} the unit of description is the current state and its exits, not a catalog of everything the server could ever do.
    \item \textbf{Graceful fallback:} when the agent fails, the human can operate the same page directly. Early evidence that machine-legible interfaces improve agent reliability is consistent with this picture \cite{schultze2025building}.
\end{enumerate}

\subsection{Grounded Human--Agent Interaction}\label{sec:bg:hai}

% 

% Limitations, <= supports 

\subsection{Summary}\label{sec:bg:gap}

Web agents can operate human-facing sites, but they infer their actions from artifacts made for human perception, and the inference is brittle. Generative-UI frameworks can render dynamic surfaces, but they ground the actions those surfaces expose weakly: inferred from the page, fixed at build time, or delegated to a machine-only contract that runs parallel to the human interface and can drift from it. No existing framework binds a human-facing generated surface to a state-dependent, server-authoritative action set.

This thesis responds by using HTML hypermedia affordances as that binding layer: the generated interface is constrained to the controls the server placed in the current representation, so every action it offers is currently valid. This is also the precise sense in which the title claims determinism. What is deterministic is the action space, not the agent: the set of offerable actions is a function of server state alone rather than of model sampling; every offered control exists and is executable at the moment it is shown; every refusal is structural. The agent's selection among valid actions remains stochastic, and no stronger claim is made. The following chapter defines this property precisely and builds the framework that realizes it. 


\clearpage

\section{Design}
\label{sec:design}
The previous chapter ended with a gap: no existing systems binds a human-facing generated surface to an action set that is both state-dependent and owned by the server. This chapter defines that binding.

\subsection{Overview}\label{sec:design:overview}

The framework consists of one idea applied at two places. A user states an intent. The agent fetches the application's current representation and extracts from it the controls the server has placed there. Those controls, taken together, form a closed set --- the \emph{affordance set} of the current state. 

Everything downstream is confined to that set: the agent may only execute actions drawn from it, and the interface presented to the user may only offer controls drawn from it. Because the set is derived from the representation the server chose to serve, and because a well-behaved server renders exactly the controls that are valid in the current state, an action outside the set is unavailable to both the agent and the interface.


\subsection{The Affordance Set}\label{sec:design:formal}

\todo{figure: server state $s \rightarrow$ representation $r(s) \rightarrow$ affordance set $A(s)$, then two branches executing and presenting}

\subsubsection{Definitions}\label{sec:design:defs}

Let $s$ denote the application state held by the server, and $r(s)$ the representation the server returns for that state. A \emph{parser} extracts from a representation the controls declared in its vocabulary; for HTML those are links, forms, and the attributes of server-driven extensions such as htmx \cite{htmxlDocs}. The \emph{affordance set} is the result:

\[ 
A(s) \;=\; \mathrm{parse}\bigl(r(s)\bigr).
\] 

Each element of $A(s)$ is a control reduced to what is needed to exercise it: a method, a target, an ordered list of declared input fields, and a human-readable label. A field carries its name, its type, whether it is required, any value fixed by the control, and, where the control declares them, its permitted values.
% maybe delete, only use the table

\todo{table: the affordance tuple --- each component, its HTML counterpart, and why it is needed to reconstruct a request.}

$A$ is indexed by state: it describes what can be done \emph{now}, not what the application can ever do. 


\subsubsection{The soundness}\label{sec:design:soundness}

Every consumer of the affordance set works by selecting from it. It may select one element to execute, or a subset to present, but it does not author elements of its own. The consequence is trivial:

\begin{quote}
\emph{Soundness.} Anything a consumer executes or presents is an element of $A(s)$, and is therefore valid in $s$.
\end{quote}

An invalid action lies outside the space being selected from. This is the precise sense in which the title of this thesis claims determinism: what is deterministic is the action space, being a function of server state alone, not the agent's choice within it.


These are not claimed...

\begin{enumerate}
    \item \textbf{Reachability is bounded by disclosure.} Actions the application supports but does not expose in its representation are unreachable through this construction. This is a genuine cost, and Chapter~\ref{sec:evaluation} measures it.
    \item \textbf{Selection remains stochastic.} Which element of $A(s)$ gets chosen is a model decision with all the usual uncertainty. Only the set is guaranteed. Selecting and ordering what is relevant to the user's intent is the agent's job.
\end{enumerate}

\subsection{Constructing $A(s)$}

\todo{need more test before writing}

\todo{existing HTML parsers like ASTs}

\subsubsection{The parser}\label{sec:design:parser}

\subsubsection{Loss}\label{sec:design:parser-loss}



\subsection{Contracts}\label{sec:design:contracts}

Soundness holds only if both ends of the exchange behave. 

\subsubsection{The site contract}\label{sec:design:site-contract}

\begin{enumerate}
    \item[\textbf{S1}] \textbf{Only valid controls are rendered.} The representation for $s$ contains a control if and only if the corresponding action is valid in $s$.
    \item[\textbf{S2}] \textbf{Requests are reconstructible.} Every value a control submits appears as a declared field. No value is computed at request time by client-side script. %!! This is tough.
    \item[\textbf{S3}] \textbf{The server guards its own state.} A request for an action invalid in the current state is refused.
\end{enumerate}

S1 is the substantive one and the framework's main adoption cost: it asks applications to treat their representation as an honest report of what is currently possible. Many applications already do this, since it is also what makes a page usable by a person; but many do not, and Chapter~\ref{sec:discussion} returns to what this means in practice.

\todo{cite: find some examples}

\subsubsection{The client contract}\label{sec:design:client-contract}

% may move to implementation
\begin{enumerate}
    \item[\textbf{C1}] \textbf{No site knowledge.} The client holds no site-specific identifiers, urls, or business vocabulary. A control is identified only by what the representation declares about it.
    %!! this actually leads to problem with many... like cookies or other things that needed to be passed, may discuss or what
    \item[\textbf{C2}] \textbf{Vocabulary extensibility.} Supporting a further hypermedia format means adding a parser that produces the same affordance shape.
\end{enumerate}

\subsection{The Consumers of \texorpdfstring{$A(s)$}{A(s)}}\label{sec:design:consumers}

\subsubsection{Execution}\label{sec:design:consumer-agent}

At the execution boundary the agent does three things with the affordance set: it selects an element, supplies values for those declared fields it can derive from the user's intent, and issues the corresponding request. Values it cannot derive are not invented; they are asked of the user, which leads directly to the next section.

\subsubsection{Presentation}\label{sec:design:consumer-ui}

\paragraph{The action layer.} The controls offered to the user. This layer is bounded: $\mathrm{controls}(UI) \subseteq A(s)$. Each control can be rendered directly from an affordance, one generic renderer per declared field type, with no model in the rendering path and no task-specific component registry. 
%!! or ... not quite sure about rendering

\paragraph{The elicitation layer.} What the interface asks the user for. This layer is also grounded in $A(s)$, but through the fields of its elements. Which questions must be asked is determined by which required fields cannot be derived from the intent; how they are asked is determined by the declared field types; and what counts as an admissible answer is determined by the permitted values a control declares.

\paragraph{The presentation layer.} Everything else: the data shown, the grouping, the wording, the explanation of what is offered and of what is structurally absent. This layer draws on the non-interactive content of $r(s)$ and on the agent's own language, and it is not bounded by anything in this framework. The division can be summarised as follows: \emph{the content is narrated by the agent; the controls are authorised by the server.}

\todo{figure: one affordance $\rightarrow$ one rendered control, field by field (name / type / required / permitted values / fixed value $\rightarrow$ widget), showing that the rendering step needs no model.}

\todo{table: field type $\rightarrow$ rendered control $\rightarrow$ what the user may change $\rightarrow$ what the control fixes.}

\subsection{Composition and Scope}\label{sec:design:composition}

Nothing above depends on there being only one application. If a surface draws on sources $1 \dots n$, each in its own state, the same reasoning gives
\[ \mathrm{controls}(UI) \;\subseteq\; \bigcup_{i} A_i(s_i),
\] so soundness is preserved under composition, provided each source supplies a state-dependent valid set.

\todo{table: extend the framework comparison of Section~\ref{sec:bg:frameworks} with the two axes introduced here: (i) is the action set state-dependent? (ii) is it the same artifact as the human interface? Rows: DOM-inference agents, build-time registries, tool catalogs such as MCP, this framework. Reference the table from both chapters.}

The framework applies to applications that satisfy the site contract, and it is therefore as much a proposal to application authors as a client architecture. Within such applications, reachability is bounded by what the server discloses; interface quality, dialogue design, and free-form information gathering fall outside the property; and composition across several sources is defined here but only exercised for a single source in what follows.


\clearpage

\section{Implementation}
\label{sec:implementation}
This chapter describes the artifact built to make the framework of Chapter~\ref{sec:design} concrete and testable. It has two independently runnable parts: a fixed generic client and one or more hypermedia sites. 

\todo{figure: architecture. User $\leftrightarrow$ client (discover / choose / render / replay) $\leftrightarrow$ site over HTTP.}

\subsection{The Site Side}\label{sec:impl:site}

The first site, \texttt{b1}, is a minimal but fully browsable restaurant-reservation application...

\todo{Introduce b1,}

\subsubsection{The state machine as the authority}\label{sec:impl:site:model}

\todo{we have a state machine that can be used as ground truth. states: browsing, holding, details entered, confirmed, conflict, cancelled. Note the second role: it is the authoritative definition of validity, hence the annotation-free ground truth used in Chapter~\ref{sec:evaluation}.}

\todo{table: the six states and the affordances exposed in each, highlighting the deliberately absent controls.}

\todo{figure: state diagram, transitions labelled by the affordance that causes them.}

\subsubsection{A representation layer that cannot disobey}\label{sec:impl:site:render}

\todo{The renderer is  dumb: every interactive control it emits comes from the valid-affordance function and it has no other source.}

\subsubsection{Vocabulary}\label{sec:impl:site:vocab}

\todo{Core flow (browse, hold, enter details, confirm) carried by standard HTML: \texttt{<a href>} for navigation, \texttt{<form action method>} with named inputs for operations. htmx \citep{htmxlDocs} used only at the two boundaries native HTML cannot describe about itself --- honest verbs beyond GET/POST (release, cancel via \texttt{hx-delete}) and partial refresh (availability via \texttt{hx-get} with \texttt{hx-target}). Both readable from the returned HTML with no JavaScript executed.}


\subsection{Discovery}\label{sec:impl:discovery}

\subsubsection{The normalised affordance}\label{sec:impl:discovery:shape}

\todo{One shape for every control, native or htmx: method, url as written, ordered declared fields, human-visible label, source vocabulary. A field records name, type, required, any prefilled or fixed value, and the enumerated options of a select. Identity for de-duplication includes each field's name and value: five ``hold this slot'' forms posting to the same url with different hidden slot values are five distinct affordances and would wrongly collapse if values were ignored.}

\todo{figure or listing: one rendered control $\rightarrow$ its normalised tuple, paired with the Design figure so the reader sees representation $\rightarrow$ tuple $\rightarrow$ re-rendered control as one round trip.}

\subsubsection{Vocabulary parsers}\label{sec:impl:discovery:parsers}

\todo{Native parser: \texttt{<a href>} as GET navigation, \texttt{<form>} as operation, but only when the form has a real submit control, so an action-less form whose only button carries htmx attributes yields no phantom affordance. htmx parser: \texttt{hx-get/post/put/delete} give method and url; where a control declares that it includes the enclosing form, that form's inputs become its fields. Presentation-only attributes (\texttt{hx-target}, \texttt{hx-swap}) ignored, since the client always reads whole responses. }

\subsubsection{Closure}\label{sec:impl:discovery:closure}

\todo{Discovery parses once, runs both parsers, de-duplicates. The resulting list \emph{is} $A(s)$, complete before any consumer sees it}

\subsection{Execution}\label{sec:impl:agent}

\subsubsection{Choosing}\label{sec:impl:agent:choose}

\todo{The only place nondeterminism lives: one interface handed the goal and the already-closed set, returning an index into it or a stop signal. Interchangeable backends --- manual, an OpenAI-compatible chat endpoint (by default a small local model via Ollama, temperature 0, strict JSON, one reprompt, fallback to manual), and a hosted frontier model. The prompt instructs the model to choose a single action and supply values only for fields derivable from the goal, never inventing personal data; missing required fields are asked of the user. Note this backend list is also the model ladder of Chapter~\ref{sec:evaluation}, which adds a random backend.}

\subsubsection{Acting by replay}\label{sec:impl:agent:replay}

\todo{construct exactly the request the control describes --- query parameters for GET, form-encoded body otherwise, fixed fields forced back to their declared value because they are determined by the control rather than chosen. The client never sends the header marking a partial request, so the site answers on its ordinary Post/Redirect/Get path and the redirect is followed; the client therefore always reads a complete next representation, with no fragments and no client-side runtime to interpret. ``Clicking the control'' is thus a pure function of method, url, fields and values.}

\subsubsection{The loop}\label{sec:impl:agent:loop}

\todo{Cold-start from an entry url; at each step read the representation, discover, choose, fill any required field the chooser could not supply, replay, read the next representation. No map of states or urls is carried; each state is learned only by fetching it. This out-of-band-free progressive discovery is what makes ``hypermedia-constrained'' literal.}

\todo{figure: sequence diagram of one iteration --- GET, discover, choose, fill, replay, GET next.}

\subsection{Presentation}\label{sec:impl:ui}

%!!NOT YET BUILT. 
\todo{The current artifact exposes the affordance set only through a command-line interface, which covers the execution boundary but leaves the rendering boundary of Section~\ref{sec:design:consumer-ui} unimplemented. Write this section once the renderer exists.}

\todo{Planned content: (i) one generic renderer per declared field type, no model in the rendering path, so the surface is a function of the subset handed to it --- it cannot contain a control outside $A(s)$ nor ask for an input no affordance declares; (ii) what the agent contributes --- subset selection, ordering, prefilling from the goal, and explanation of both what is offered and what is structurally absent, never authoring a control; (iii) a human click and an agent step go through the same replay path, so the two operate one action set rather than two.}

\todo{figure: screenshots of the generated surface for two different states, side by side, with the corresponding affordance set printed beside each.}

\todo{Takeover and fallback: the human may act at any step, and if the agent fails the underlying site remains directly operable --- the graceful-fallback consequence of Section~\ref{sec:bg:hypermedia} realised.}

\subsection{A Second Site}\label{sec:impl:b2}

\todo{To be written once built. A second site in a different domain (the study-plan approval workflow of earlier drafts) shows that the same unchanged client operates a different application. Keep it short --- an existence claim, with evidence reported in Chapter~\ref{sec:evaluation}. State explicitly what had to change in the client, the intended answer being nothing.}

\subsection{Technology Stack and Reproducibility}\label{sec:impl:stack}

\todo{Node.js and TypeScript run directly with tsx, no build step, managed with pnpm. Sites use Express and template-literal HTML so affordances are visible in the source; the client uses platform fetch for replay and linkedom to parse responses. Shared toolchain at the top level, one TypeScript configuration for the workspace.}

\todo{Reproducibility scripts: run the site, drive its flows over HTTP and assert the refusals, self-check the client engine against a running site, and drive the site end to end from a natural-language goal. Type-check and both self-checks pass; the full happy path has been driven unattended by a small local model.}

\todo{Describe the structured log, the sole data source of Chapter~\ref{sec:evaluation}: per step, the task, condition, model, state, exposed affordance set, selected action, supplied values, resulting state. Fixing this schema is a prerequisite for the experiments.}


\clearpage

\section{Evaluation}
\label{sec:evaluation}

\subsection{What Is Evaluated}\label{sec:eval:aim}

\todo{table: the framing table. Rows precision / recall; columns unconstrained action sources / this framework; cells measured, measured, ``guaranteed, not measured'', measured. Small and early: it defines the chapter.}

\todo{State the three questions, together answering RQ3.
Q1 \emph{Magnitude}: how often does an action source without state-dependent affordances lead a system to offer actions invalid in the current state?
Q2 \emph{Cost}: what share of legitimate user goals becomes unreachable when the action space is restricted to the controls the server rendered?
Q3 \emph{Benefit}: to what extent does the constraint substitute for model capability?}

\todo{Scope exclusions: interface aesthetics, completion time, user preference, and anything requiring human participants.}

\subsection{Method}\label{sec:eval:method}

\subsubsection{Ground truth}\label{sec:eval:method:truth}

\todo{The site's valid-affordance function (Section~\ref{sec:impl:site:model}) is the authoritative definition of what is legal in a state, so every proposed action and every rendered control is labelled automatically by consulting the same function the site itself uses.}

\subsubsection{Conditions}\label{sec:eval:method:conditions}

\todo{Same client, same backend throughout; only the surface through which the application declares what can be done varies.
H \emph{hypermedia} --- the server-rendered representation, whose controls are the currently valid set.
J \emph{data only} --- the same backend as a plain JSON state resource (entity data and status, no action declarations); stands for a model placed over a conventional data API and asked to produce an interface.
%{ "booking": { "status": "confirmed", "time": "19:00", "guests": 2 } }
J+ \emph{static catalog} --- the same JSON state plus a complete, correctly typed catalog of every operation the application supports, unfiltered by state; stands for build-time tool registries and server-declared tool catalogs such as MCP.}
% hold_slot(slot), enter_details(name, phone), confirm(), cancel(), change_slot(slot)

\todo{J+ is a steelman. It gives the baseline an action vocabulary that is machine-readable, typed and server-authored, so the only thing it lacks relative to H is state-dependence.}

\todo{Fairness controls: same model, prompt scaffolding, goal text and step budget across conditions, with all three prompts in an appendix. What necessarily differs is the shape of the action information, which is the intended manipulation.}

\todo{table: the three conditions against state information, action information, state-dependent, same artifact as the human interface, and real-world analogue.}

\todo{figure: experiment design --- one backend, three surfaces, one client, four chooser backends, one log schema.}


\subsubsection{Task set}\label{sec:eval:method:tasks}

\todo{Goals derived from the state machine rather than chosen by hand, so the task set cannot be suspected of having been selected to favour the framework; describe the derivation procedure.}

\begin{enumerate}
    \item T1 \emph{valid and exposed} --- legal in the current state and present in the representation; every condition can in principle succeed.
    \item T2 \emph{valid but not exposed} --- legitimate in the domain but not offered in the current representation. H must fail, refusing something in fact possible, whereas an unconstrained condition may guess correctly. This class measures the cost admitted in Section~\ref{sec:design:soundness}.
    \item T3 \emph{invalid in the current state} --- editing a confirmed booking, forcing a submission while in conflict. Desired behaviour is a correct and explained refusal; the failure mode is offering a control that does not exist.
\end{enumerate}

\todo{table: the task set --- identifier, class, goal text, starting state, target state, expected outcome under H. Full listing in an appendix if long.}

\subsubsection{Models and repetition}\label{sec:eval:method:models}

\todo{4 chooser backends spanning capability: a hosted frontier model, a mid-tier model, a small local model, and a degenerate backend selecting uniformly at random. Justify the random backend --- under H it still selects from a set of valid actions, so whatever it achieves is achieved by the constraint rather than the model, making it the cleanest single piece of evidence for Q3.}

\subsubsection{Metrics}\label{sec:eval:method:metrics}

\todo{Define each in terms of log fields, with its expected value under H.}

\begin{enumerate}
    \item invalid action rate
    \item invalid control rate
    \item task success
    \item correct refusal rate
    \item false refusal rate
    \item steps to completion
    \item offer consistency
    \item unfounded elicitation rate
    \item ...
\end{enumerate}

\todo{table: metric, definition, log fields, question served, expected value under H.}

\subsection{Results}\label{sec:eval:results}

\todo{Pending the experiments.}

\subsubsection{A worked trace}\label{sec:eval:results:trace}

\todo{Before any number, one narrated end-to-end run: the goal, the states traversed, the affordance set at each step, the action selected, the surface generated. Include one refusal trace from T3}

\subsubsection{Q1: the size of the problem}\label{sec:eval:results:q1}


\subsubsection{Q2: what the constraint costs}\label{sec:eval:results:q2}


\subsubsection{Q3: constraint against capability}\label{sec:eval:results:q3}


\subsubsection{Consistency}\label{sec:eval:results:consistency}


\subsubsection{Cross-site generalisation}\label{sec:eval:results:b2}


\subsection{Threats to Validity}\label{sec:eval:threats}

\todo{To be written with the results}

\todo{Closing paragraph: map results back to RQ3 and forward to Chapter~\ref{sec:discussion} for what they imply for application authors.}


\clearpage

\section{Discussion}
\label{sec:discussion}
The idea of discussion is to bring up points that you noticed during research.
As a chapter, this is the most freeform one.
A good way to discuss the points is first to write up subtitles with clear statements and then expand on those.
Discussion points can occasionally be also in a question format or refer to a larger context beyond the thesis to bring in ideas from other papers to consider.
In case your work has limitations that you can observe, it can be a good idea to write a separate section considering those.


\clearpage

\section{Conclusion}
\label{sec:conclusion}
The idea of the conclusion is the opposite of the introduction.
Here we literally bring the reader out of the topic, not in.
A good way to write conclusion is to start by reminding the reader of your research question(s) and then state how your work addressed them.
To finish your conclusion, consider open research questions (this can be a list for example).


\clearpage

\bibliographystyle{unsrtnat}
\bibliography{references/HAI, references/WEB, references/UI, references/MISC}


\end{document}
